% Technical note: O7 saturation parameter and Lorentz factor from projective capacity.
% Status: sections 1-3 established; section 4 contains the open convergence lemma
% anchored to Q5b/Q11; section 5 is an algebraic corollary.
% Dependencies: Beau2026a7 (O7), Beau2026q5b, Beau2026q11, Beau2026c (BI).

\documentclass[11pt]{article}
\usepackage{amsmath,amsthm,amssymb}
\usepackage[colorlinks=true,linkcolor=blue,citecolor=blue]{hyperref}
\usepackage{doi}

\newtheorem{lemma}{Lemma}
\newtheorem{theorem}{Theorem}
\newtheorem{corollary}{Corollary}
\newtheorem{remark}{Remark}

\title{Spectral Mobility and the Lorentz Factor\\
from Projective Capacity Saturation}
\author{J\'er\^ome Beau}
\date{2026}

\begin{document}
  \maketitle

  \begin{abstract}
    The O-series saturation parameter $\eta_n^{\mathrm{O7}}$ is defined in~\cite{Beau2026a7}
    as the ratio of cumulative occupied projective capacity to residual frontier novelty.
    We show that this ratio is not an arbitrary fitting variable: under the
    Born--Infeld admissibility budget established in~\cite{Beau2026c}, a quadratic
    normalization of the two capacity components implies
    $\eta_n^{\mathrm{O7}} \to \beta\gamma$,
    where $\beta\gamma$ is the relativistic mobility and $\gamma$ is the Lorentz factor.
    Consequently, the O-series profile $\Phi(\eta) = 1/\sqrt{1+\eta^2}$ converges to
    $1/\gamma$: the inverse Lorentz factor is the residual-capacity profile of the
    admissible spectral cascade.
    The quadratic normalization is not postulated; it is expected to follow from the
    Carnot--Carath\'eodory continuum geometry and co-metric closure established
    in~\cite{Beau2026q5b,Beau2026q11}.
    The main convergence lemma, identifying the O7 observables as the normalized
    components of the Born--Infeld budget, is stated precisely and left as the
    remaining technical task.
  \end{abstract}

  \tableofcontents
  \newpage

  \section{O7 occupation and residual novelty}
    \label{sec:o7}

    The O7 saturation parameter is defined in~\cite{Beau2026a7} by
    \[
      \eta_n^{\mathrm{O7}}(x)
      =
      \frac{B_n(x)}{\Sigma_n(x)},
      \qquad
      B_n(x) := |\psi_n(x)|^2
      =
      \frac{|S_n \cap C_x|}{|C_x|}.
    \]
    Here $C_x$ is the effective resolution cell associated with the observable state $x$,
    and $S_n$ is the BFS-reachable set up to cascade depth $n$.

    The observable $B_n(x)$ is cumulative.
    It measures the fraction of the cell $C_x$ already reached by the BFS cascade by
    depth $n$.
    In particular,
    \[
      0 \leq B_n(x) \leq 1,
      \qquad
      B_{n+1}(x) \geq B_n(x).
    \]
    It is therefore an occupation observable, not an infinitesimal mobility.

    By contrast, $\Sigma_n(x)$ measures the average effective novelty contributed by
    the BFS frontier at depth $n$: the mean fraction of genuinely new projective
    direction brought by admissible paths reaching $\partial S_n$ within $C_x$.
    It is a residual frontier observable, instantaneous and non-increasing in $n$.

    The ratio $\eta_n^{\mathrm{O7}}$ is therefore not a fitting parameter.
    It is an occupied-over-residual capacity ratio:
    \[
      \eta_n^{\mathrm{O7}}
      =
      \frac{\text{cumulative occupied capacity}}{\text{residual frontier novelty}}.
    \]

    The question addressed in this note is whether this O7 ratio converges, in the
    continuum geometric limit, to the relativistic mobility ratio $dI_{\mathrm{space}}/dI_\tau$.

  \section{Algebraic Born--Infeld target}
    \label{sec:bi}

    The cascade depth $n$ is not a time coordinate.
    It is the ordinal depth of an admissible projective history, counting the number of
    distinguishable projective transitions undergone along that history~\cite{Beau2026c}.

    Assume a bounded Born--Infeld mobility budget at each projective depth
    increment $dn$:
    \[
      dI_\tau^2 + dI_{\mathrm{space}}^2 = c_{\mathrm{BI}}^2\,dn^2.
    \]
    Define the normalized spectral mobility
    \[
      \beta := \frac{dI_{\mathrm{space}}}{c_{\mathrm{BI}}\,dn}.
    \]
    Then the residual temporal component is
    \[
      dI_\tau = c_{\mathrm{BI}}\sqrt{1-\beta^2}\,dn,
    \]
    and therefore
    \[
      \frac{dI_{\mathrm{space}}}{dI_\tau}
      =
      \frac{\beta}{\sqrt{1-\beta^2}}
      =
      \beta\gamma,
      \qquad
      \gamma := \frac{1}{\sqrt{1-\beta^2}}.
    \]
    If a spectral saturation parameter $\eta$ satisfies $\eta = dI_{\mathrm{space}}/dI_\tau$, then
    \[
      \Phi(\eta) = \frac{1}{\sqrt{1+\eta^2}} = \frac{1}{\gamma}.
    \]
    The inverse Lorentz factor is the residual-capacity profile of a quadratic
    Born--Infeld budget.
    This step is purely algebraic.
    The non-trivial question is whether the O7 observables $B_n$ and $\Sigma_n$
    converge to the two normalized components of this budget.

  \section{Why the reduced filling model is insufficient}
    \label{sec:linear}

    The reduced filling model of~\cite{Beau2026a7} gives, before geometric correction,
    \[
      B_n \simeq \frac{n}{D},
      \qquad
      \Sigma_n \simeq 1 - \frac{n}{D},
    \]
    equivalently $B_n + \Sigma_n \simeq 1$.
    This is a linear residual law.
    It describes monotone filling of a finite resolution cell but does not produce
    the Lorentzian residual factor.
    Indeed, setting $B_n \to \beta$ in the linear model would give
    $\Sigma_n \to 1-\beta$,
    whereas the Born--Infeld mobility budget requires $\Sigma_n \to \sqrt{1-\beta^2}$.

    The linear filling model therefore captures the monotone occupation mechanism
    but not the continuum geometric correction responsible for the Lorentzian
    square root.
    The required structure is the quadratic residual law
    \[
      B_n^2 + \Sigma_n^2 \to 1.
    \]
    Under the identification $B_n \to \beta$, this gives $\Sigma_n \to \sqrt{1-\beta^2}$.

  \section{From linear filling to quadratic normalization via Q5b}
    \label{sec:q5b}

    \subsection{Q5b input}

      The reduced filling model operates at the combinatorial level of a finite
      resolution cell.
      The Q5b continuum limit supplies the missing geometric structure~\cite{Beau2026q5b}.

      In the pre-saturation large-$q$ regime, BFS shells on
      $\operatorname{Heis}_3(\mathbb{Z}/q\mathbb{Z})$ converge to
      Carnot--Carath\'eodory balls on $\operatorname{Heis}_3(\mathbb{R})$.
      The linear residual law of the reduced model is a coarse filling approximation
      valid at finite $D$ and finite $q$; it does not encode the sub-Riemannian
      anisotropy of the CC geometry.

      The effective co-metric extracted from the completed principal-symbol geometry
      of the admissibility operator is Lorentzian in the homogeneous regime.
      After the Q11 co-metric closure~\cite{Beau2026q11}, the normalized co-metric is
      \[
        g^{\mu\nu} = 2\,\eta^{\mu\nu}, \qquad \eta = \mathrm{diag}(-1,+1,+1,+1),
      \]
      with no free parameter.
      This Lorentzian structure imposes a quadratic normalization on the effective
      admissible components in the continuum limit — replacing the linear complement
      of the reduced filling model.

      The chain is therefore:
      \begin{align*}
        &\text{CC convergence (Q5b)} &&\Rightarrow&&
        \text{continuum spatial and temporal components exist,}\\
        &\text{principal symbol / Q11 closure} &&\Rightarrow&&
        \text{quadratic normalization of effective components,}\\
        &\text{observable identification (Lemma~\ref{lem:main})} &&\Rightarrow&&
        B_n^2 + \Sigma_n^2 \to 1.
      \end{align*}
      The geometric structure is not postulated; it comes from Q5b and Q11.
      The remaining task is the third step: showing that the O7 observables $B_n$
      and $\Sigma_n$ are precisely the normalized projections of the spatially occupied
      and temporally residual components of the continuum budget.

    \subsection{Main convergence lemma}

      \begin{lemma}[O7 observables as normalized CC budget components]
        \label{lem:main}
        Let $B_n(x) := |\psi_n(x)|^2 = |S_n \cap C_x|/|C_x|$ be the O7 cumulative
        occupation observable, and let $\Sigma_n(x)$ be the O7 residual frontier
        novelty~\cite{Beau2026a7}.
        In the pre-saturation large-$q$ regime, under the Q5b BFS-to-CC convergence
        and the Q11 homogeneous co-metric closure, these observables converge to the
        normalized spatial and temporal components of the Born--Infeld mobility budget:
        \[
          B_n \longrightarrow \beta,
          \qquad
          \Sigma_n \longrightarrow \sqrt{1-\beta^2},
        \]
        where $\beta = v_{\mathrm{eff}}/c_{\mathrm{BI}}$ is the normalized projected spectral mobility.
        Equivalently,
        \[
          B_n^2 + \Sigma_n^2 \longrightarrow 1.
        \]
      \end{lemma}

      The remaining work is the observable-identification step.
      Q5b provides the continuum quadratic geometry.
      O7 provides the occupied-over-residual observables.
      The lemma asserts that these two descriptions refer to the same normalized
      capacity components in the common large-$q$, pre-saturation regime.

  \section{Consequence for the Lorentz factor}
    \label{sec:consequence}

    \begin{theorem}[Lorentz factor as residual projective capacity]
      \label{thm:lorentz}
      Assume the Born--Infeld mobility budget of Section~\ref{sec:bi} and the
      convergence of Lemma~\ref{lem:main}.
      Then the O7 saturation parameter satisfies
      \[
        \eta_n^{\mathrm{O7}}
        =
        \frac{B_n}{\Sigma_n}
        \longrightarrow
        \frac{\beta}{\sqrt{1-\beta^2}}
        =
        \beta\gamma,
      \]
      and therefore the O-series saturation profile satisfies
      \[
        \Phi\!\left(\eta_n^{\mathrm{O7}}\right)
        =
        \frac{1}{\sqrt{1+\left(\eta_n^{\mathrm{O7}}\right)^2}}
        \longrightarrow
        \frac{1}{\gamma}.
      \]
    \end{theorem}

    The special-relativistic residual factor is thus recovered from the O7
    occupied-over-residual capacity ratio, provided the Q5b continuum geometry
    upgrades the reduced linear filling law to the quadratic Born--Infeld residual
    law.
    The Lorentz factor is not added as an external kinematical postulate.
    It is the effective reading of residual temporal capacity in a bounded
    projective budget.

    \begin{remark}[Saturation pressure at $v \to c$]
      Under the identification $\eta_n^{\mathrm{O7}} \to \beta\gamma$, the O7 saturation
      pressure $P_n^{\mathrm{sat}} = \gamma\,\eta_n$ satisfies
      $P_n^{\mathrm{sat}} \to \beta\gamma^2$, which diverges as $v \to c$.
      This is consistent with the light cone reading $dI_\tau = 0 \Leftrightarrow v = c$:
      the budget is fully consumed by spatial mobility and no temporal residual
      remains.
    \end{remark}

    \begin{remark}[Scope]
      Theorem~\ref{thm:lorentz} directly yields time dilation
      ($d\tau/dn = 1/\gamma$) and the light cone ($dI_\tau = 0 \Leftrightarrow v = c$).
      Length contraction requires the additional reconstruction of effective Lorentz
      transformations between inertial observers.
      The extension to general relativity follows from the non-uniform regime: when
      localised configurations modify the local spectral budget, the position-dependent
      metric $g_{\mu\nu}(x)$ encodes gravitational time dilation as a comparison of
      residual temporal capacities along distinct projected histories.
    \end{remark}

    \bibliographystyle{plain}
    \bibliography{cosmochrony-bibliography}

\end{document}
